% Options for packages loaded elsewhere
\PassOptionsToPackage{unicode}{hyperref}
\PassOptionsToPackage{hyphens}{url}
\documentclass[
  12pt,
]{ctexart}
\usepackage{xcolor}
\usepackage{amsmath,amssymb}
\setcounter{secnumdepth}{-\maxdimen} % remove section numbering
\usepackage{iftex}
\ifPDFTeX
  \usepackage[T1]{fontenc}
  \usepackage[utf8]{inputenc}
  \usepackage{textcomp} % provide euro and other symbols
\else % if luatex or xetex
  \usepackage{unicode-math} % this also loads fontspec
  \defaultfontfeatures{Scale=MatchLowercase}
  \defaultfontfeatures[\rmfamily]{Ligatures=TeX,Scale=1}
\fi
\usepackage{lmodern}
\ifPDFTeX\else
  % xetex/luatex font selection
\fi
% Use upquote if available, for straight quotes in verbatim environments
\IfFileExists{upquote.sty}{\usepackage{upquote}}{}
\IfFileExists{microtype.sty}{% use microtype if available
  \usepackage[]{microtype}
  \UseMicrotypeSet[protrusion]{basicmath} % disable protrusion for tt fonts
}{}
\makeatletter
\@ifundefined{KOMAClassName}{% if non-KOMA class
  \IfFileExists{parskip.sty}{%
    \usepackage{parskip}
  }{% else
    \setlength{\parindent}{0pt}
    \setlength{\parskip}{6pt plus 2pt minus 1pt}}
}{% if KOMA class
  \KOMAoptions{parskip=half}}
\makeatother
\usepackage{color}
\usepackage{fancyvrb}
\newcommand{\VerbBar}{|}
\newcommand{\VERB}{\Verb[commandchars=\\\{\}]}
\DefineVerbatimEnvironment{Highlighting}{Verbatim}{commandchars=\\\{\}}
% Add ',fontsize=\small' for more characters per line
\newenvironment{Shaded}{}{}
\newcommand{\AlertTok}[1]{\textcolor[rgb]{1.00,0.00,0.00}{\textbf{#1}}}
\newcommand{\AnnotationTok}[1]{\textcolor[rgb]{0.38,0.63,0.69}{\textbf{\textit{#1}}}}
\newcommand{\AttributeTok}[1]{\textcolor[rgb]{0.49,0.56,0.16}{#1}}
\newcommand{\BaseNTok}[1]{\textcolor[rgb]{0.25,0.63,0.44}{#1}}
\newcommand{\BuiltInTok}[1]{\textcolor[rgb]{0.00,0.50,0.00}{#1}}
\newcommand{\CharTok}[1]{\textcolor[rgb]{0.25,0.44,0.63}{#1}}
\newcommand{\CommentTok}[1]{\textcolor[rgb]{0.38,0.63,0.69}{\textit{#1}}}
\newcommand{\CommentVarTok}[1]{\textcolor[rgb]{0.38,0.63,0.69}{\textbf{\textit{#1}}}}
\newcommand{\ConstantTok}[1]{\textcolor[rgb]{0.53,0.00,0.00}{#1}}
\newcommand{\ControlFlowTok}[1]{\textcolor[rgb]{0.00,0.44,0.13}{\textbf{#1}}}
\newcommand{\DataTypeTok}[1]{\textcolor[rgb]{0.56,0.13,0.00}{#1}}
\newcommand{\DecValTok}[1]{\textcolor[rgb]{0.25,0.63,0.44}{#1}}
\newcommand{\DocumentationTok}[1]{\textcolor[rgb]{0.73,0.13,0.13}{\textit{#1}}}
\newcommand{\ErrorTok}[1]{\textcolor[rgb]{1.00,0.00,0.00}{\textbf{#1}}}
\newcommand{\ExtensionTok}[1]{#1}
\newcommand{\FloatTok}[1]{\textcolor[rgb]{0.25,0.63,0.44}{#1}}
\newcommand{\FunctionTok}[1]{\textcolor[rgb]{0.02,0.16,0.49}{#1}}
\newcommand{\ImportTok}[1]{\textcolor[rgb]{0.00,0.50,0.00}{\textbf{#1}}}
\newcommand{\InformationTok}[1]{\textcolor[rgb]{0.38,0.63,0.69}{\textbf{\textit{#1}}}}
\newcommand{\KeywordTok}[1]{\textcolor[rgb]{0.00,0.44,0.13}{\textbf{#1}}}
\newcommand{\NormalTok}[1]{#1}
\newcommand{\OperatorTok}[1]{\textcolor[rgb]{0.40,0.40,0.40}{#1}}
\newcommand{\OtherTok}[1]{\textcolor[rgb]{0.00,0.44,0.13}{#1}}
\newcommand{\PreprocessorTok}[1]{\textcolor[rgb]{0.74,0.48,0.00}{#1}}
\newcommand{\RegionMarkerTok}[1]{#1}
\newcommand{\SpecialCharTok}[1]{\textcolor[rgb]{0.25,0.44,0.63}{#1}}
\newcommand{\SpecialStringTok}[1]{\textcolor[rgb]{0.73,0.40,0.53}{#1}}
\newcommand{\StringTok}[1]{\textcolor[rgb]{0.25,0.44,0.63}{#1}}
\newcommand{\VariableTok}[1]{\textcolor[rgb]{0.10,0.09,0.49}{#1}}
\newcommand{\VerbatimStringTok}[1]{\textcolor[rgb]{0.25,0.44,0.63}{#1}}
\newcommand{\WarningTok}[1]{\textcolor[rgb]{0.38,0.63,0.69}{\textbf{\textit{#1}}}}
\usepackage{longtable,booktabs,array}
\usepackage{calc} % for calculating minipage widths
% Correct order of tables after \paragraph or \subparagraph
\usepackage{etoolbox}
\makeatletter
\patchcmd\longtable{\par}{\if@noskipsec\mbox{}\fi\par}{}{}
\makeatother
% Allow footnotes in longtable head/foot
\IfFileExists{footnotehyper.sty}{\usepackage{footnotehyper}}{\usepackage{footnote}}
\makesavenoteenv{longtable}
\setlength{\emergencystretch}{3em} % prevent overfull lines
\providecommand{\tightlist}{%
  \setlength{\itemsep}{0pt}\setlength{\parskip}{0pt}}
\usepackage{bookmark}
\IfFileExists{xurl.sty}{\usepackage{xurl}}{} % add URL line breaks if available
\urlstyle{same}
\hypersetup{
  hidelinks,
  pdfcreator={LaTeX via pandoc}}

\author{SCX}
\date{2026}

\begin{document}

\section{SCX
架构清单：基底、脚本、接口、数学公式、应用场景}\label{scx-ux67b6ux6784ux6e05ux5355ux57faux5e95ux811aux672cux63a5ux53e3ux6570ux5b66ux516cux5f0fux5e94ux7528ux573aux666f}

\begin{quote}
日期：2026-06-27 \textbar{} SCX v0.4.0-pre
\end{quote}

\begin{center}\rule{0.5\linewidth}{0.5pt}\end{center}

\subsection{0.
回答你的核心问题}\label{ux56deux7b54ux4f60ux7684ux6838ux5fc3ux95eeux9898}

\textbf{Q: 做势函数时用 ACE，做
LLM/图像时用别的神经网络------怎么控制？}

A: 通过 \textbf{Encoder 抽象层}控制。SCX
不关心你的神经网络是什么结构，只要求你实现一个 \texttt{SCXStateEncoder}
接口（3 个方法：\texttt{encode}, \texttt{distance},
\texttt{cluster}）。MLIP 用 \texttt{MLIPEncoder}（手工描述符），图像用
\texttt{VisionEncoder}（CNN backbone），表格用
\texttt{TabularEncoder}。LLM 你写一个 \texttt{LLMEncoder} 把文本
embedding 成向量就行。

\textbf{Q: 基底和应用场景的脚本由接口链接？基底一样吗？}

A: \textbf{基底完全一样。} \texttt{SCXStateEncoder} 和
\texttt{SCXExpert} 是两个抽象基类，是唯一的领域接口。Encoder
以上的一切（StateDiscovery、DataClassifier、StateValue、ActionPolicy）不感知你是什么领域------它们只看到向量。换领域
= 换 Encoder 实现 + 换 YAML 配置。

\begin{center}\rule{0.5\linewidth}{0.5pt}\end{center}

\subsection{1. 总体架构图}\label{ux603bux4f53ux67b6ux6784ux56fe}

\begin{verbatim}
┌─────────────────────────────────────────────────────────────────────┐
│                        SCX v4.0 通用模块化架构                        │
├─────────────────────────────────────────────────────────────────────┤
│                                                                      │
│  应用场景层 (domains/*.yaml)                                          │
│  ┌──────────┐  ┌──────────┐  ┌──────────┐  ┌──────────┐            │
│  │ mlip.yaml│  │cifar.yaml│  │health.yaml│  │ (你的新域)│            │
│  └────┬─────┘  └────┬─────┘  └────┬─────┘  └────┬─────┘            │
│       │              │              │              │                  │
│  ─ ─ ─┼──────┬───────┼──────────────┼──────────────┼── 接口层       │
│       │      │       │              │              │                  │
│  ┌────▼──────▼───────▼──────────────▼──────────────▼────┐           │
│  │              SCXStateEncoder (抽象基类)               │           │
│  │              SCXExpert (抽象基类)                     │           │
│  └──────────────────────┬───────────────────────────────┘           │
│                          │                                            │
│  ─ ─ ─ ─ ─ ─ ─ ─ ─ ─ ─ ┼ ─ ─ ─ ─ ─ ─ ─ ─ ─ ─ ─ ─ ─ ─ 基底       │
│                          │                                            │
│  ┌──────────────────────▼───────────────────────────────┐           │
│  │  核心引擎 (不感知领域)                                 │           │
│  │  StateDiscovery → DataClassifier → StateValue         │           │
│  │  NoiseScore → RedundancyScore → LearnabilityScore     │           │
│  │  ExpertReliability → ExpertRouter → CompressStrategy  │           │
│  │  ActionPolicy → SCXFramework (orchestrator)           │           │
│  └──────────────────────────────────────────────────────┘           │
│                                                                      │
└─────────────────────────────────────────────────────────────────────┘
\end{verbatim}

\textbf{关键设计原则}：Encoder 以上不变，Encoder
以下是领域适配。新领域只需实现 Encoder + 写 YAML。

\begin{center}\rule{0.5\linewidth}{0.5pt}\end{center}

\subsection{2.
基底（核心引擎）------所有应用共享，不随领域变化}\label{ux57faux5e95ux6838ux5fc3ux5f15ux64ceux6240ux6709ux5e94ux7528ux5171ux4eabux4e0dux968fux9886ux57dfux53d8ux5316}

\subsubsection{2.1 基底的 Python
脚本清单}\label{ux57faux5e95ux7684-python-ux811aux672cux6e05ux5355}

\begin{longtable}[]{@{}
  >{\raggedright\arraybackslash}p{(\linewidth - 6\tabcolsep) * \real{0.2308}}
  >{\raggedright\arraybackslash}p{(\linewidth - 6\tabcolsep) * \real{0.3462}}
  >{\raggedright\arraybackslash}p{(\linewidth - 6\tabcolsep) * \real{0.2308}}
  >{\centering\arraybackslash}p{(\linewidth - 6\tabcolsep) * \real{0.1923}}@{}}
\toprule\noalign{}
\begin{minipage}[b]{\linewidth}\raggedright
文件
\end{minipage} & \begin{minipage}[b]{\linewidth}\raggedright
类/函数
\end{minipage} & \begin{minipage}[b]{\linewidth}\raggedright
作用
\end{minipage} & \begin{minipage}[b]{\linewidth}\centering
依赖领域吗？
\end{minipage} \\
\midrule\noalign{}
\endhead
\bottomrule\noalign{}
\endlastfoot
\texttt{core/config.py} & \texttt{SCXConfig} &
全局配置（状态数、阈值、预算等），从 JSON/YAML 加载 & ❌ \\
\texttt{core/framework.py} & \texttt{SCXFramework} & 总调度器：fit() →
状态发现 → 估值 → 分类 → 动作 & ❌ \\
\texttt{core/metrics.py} & \texttt{SCXMetrics} & 指标记录与追踪 & ❌ \\
\texttt{core/online.py} & \texttt{OnlineStateTracker},
\texttt{OnlineExpertTracker}, \texttt{OnlineSCXFramework} &
流式/在线场景的 EMA 更新 & ❌ \\
\texttt{state/discovery.py} & \texttt{StateDiscovery} &
聚类状态发现（KMeans/GMM/Spectral/HDBSCAN） & ❌ \\
\texttt{state/space.py} & \texttt{StateSpace} & 状态空间数据结构 & ❌ \\
\texttt{state/assignment.py} & \texttt{StateAssignment} &
新样本映射到已有状态 & ❌ \\
\texttt{state/metrics.py} & 状态距离/相似度度量 & 状态间距离计算 & ❌ \\
\texttt{state/robustness.py} & \texttt{StateRobustness} &
状态鲁棒性分析（misclassification/boundary confidence） & ❌ \\
\texttt{state/two\_layer.py} & \texttt{TwoLayerStateDiscovery} &
\textbf{两层描述符}：L1人工 + L2误差驱动 & ❌ \\
\texttt{valuation/classifier.py} & \texttt{DataClassifier} &
四分类器（valuable/redundant/noisy/expert\_dep） & ❌ \\
\texttt{valuation/state\_value.py} & \texttt{StateValue} + theorem-based
methods & V(s) 计算 {[}deprecated{]}, noise\_consistency\_score,
chernoff\_bound, hoeffding\_bound, feature\_strength\_diagnostic & ❌ \\
\texttt{valuation/noise\_score.py} & \texttt{NoiseScore} &
样本级/状态级噪声评分 & ❌ \\
\texttt{valuation/redundancy.py} & \texttt{RedundancyScore} & 冗余度评分
& ❌ \\
\texttt{valuation/learnability.py} & \texttt{LearnabilityScore} &
可学习性评分（一致性+非噪声） & ❌ \\
\texttt{valuation/adaptive.py} & \texttt{AdaptiveThreshold} &
自适应阈值校准（网格搜索/百分位/gap法） & ❌ \\
\texttt{valuation/influence.py} & \texttt{StateConditionedInfluence} &
状态条件影响力函数 & ❌ \\
\textbf{\texttt{yajie.py}} & \texttt{YajieCleaner} &
雅洁数据清理器（定理驱动去噪） & ❌ \\
\texttt{expert/registry.py} & \texttt{ExpertRegistry} & 专家注册与管理 &
❌ \\
\texttt{expert/reliability.py} & \texttt{ExpertReliability} &
状态条件专家可靠性 SCX\_m(s) & ❌ \\
\texttt{expert/router.py} & \texttt{ExpertRouter} &
专家路由（hard/soft） & ❌ \\
\texttt{expert/conflict.py} & 专家冲突解决 & 多专家冲突检测与仲裁 &
❌ \\
\texttt{action/compress.py} & \texttt{CompressStrategy} & 数据压缩策略 &
❌ \\
\texttt{action/acquisition.py} & 主动采集策略 & 状态条件采样 & ❌ \\
\texttt{action/policy.py} & \texttt{ActionPolicy} & 动作决策策略 & ❌ \\
\texttt{utils/data\_loader.py} & 数据加载工具 & 统一数据加载接口 & ❌ \\
\texttt{utils/evaluation.py} & 评估工具 & 指标计算、对比 & ❌ \\
\texttt{utils/visualization.py} & 可视化工具 & 状态地图、残差图、对比图
& ❌ \\
\texttt{utils/helpers.py} & 通用辅助函数 & 数组/统计工具 & ❌ \\
\end{longtable}

\textbf{总计：31 个源文件，全部领域无关。}

\subsubsection{2.2
基底使用的数学公式}\label{ux57faux5e95ux4f7fux7528ux7684ux6570ux5b66ux516cux5f0f}

\paragraph{核心定义}\label{ux6838ux5fc3ux5b9aux4e49}

\begin{longtable}[]{@{}
  >{\raggedright\arraybackslash}p{(\linewidth - 4\tabcolsep) * \real{0.2857}}
  >{\raggedright\arraybackslash}p{(\linewidth - 4\tabcolsep) * \real{0.2857}}
  >{\raggedright\arraybackslash}p{(\linewidth - 4\tabcolsep) * \real{0.4286}}@{}}
\toprule\noalign{}
\begin{minipage}[b]{\linewidth}\raggedright
公式
\end{minipage} & \begin{minipage}[b]{\linewidth}\raggedright
含义
\end{minipage} & \begin{minipage}[b]{\linewidth}\raggedright
代码位置
\end{minipage} \\
\midrule\noalign{}
\endhead
\bottomrule\noalign{}
\endlastfoot
\textbf{R\_m(s)} = E{[} ℓ(f\_m(x), f*(x)) \textbar{} x∈s {]} & 专家 m
在状态 s 的条件风险 & \texttt{expert/reliability.py} \\
\textbf{SCX\_m(s)} = P( ℓ(f\_m(x), y) \textless{} τ \textbar{} x∈s ) &
专家 m 在状态 s 下做出可接受预测的概率 &
\texttt{expert/reliability.py} \\
\textbf{V\_add(s)} = r̄(s) · ρ(s) · L(s) · {[}1−D(s){]} · max\_m
SCX\_m(s) & 采集价值（越采集越值钱） &
\texttt{valuation/state\_value.py} \\
\textbf{V\_remove(s)} = ρ(s) · (1−r̄(s)) · Sim(s) · (1−Boundary(s)) &
压缩价值（越压缩越安全） & \texttt{valuation/state\_value.py} \\
\textbf{NoiseScore(x)} = r · w\_ρ/(ρ(s)+ε) · {[}1−C(s){]} · w\_C &
噪声评分（高误差+低密度+低一致性→噪声） &
\texttt{valuation/noise\_score.py} \\
\textbf{L(s)} = C(s) · {[}1−N(s){]} & 可学习性 = 一致性 × (1−噪声) &
\texttt{valuation/learnability.py} \\
\textbf{D(s)} = f(ρ, r̄, Similarity, Boundary) & 冗余度 &
\texttt{valuation/redundancy.py} \\
\end{longtable}

\paragraph{数据四分类规则}\label{ux6570ux636eux56dbux5206ux7c7bux89c4ux5219}

\begin{verbatim}
IF expert_gap > threshold            → expert_dependent (路由)
IF r̄↑ AND ρ↑ AND C↑ AND D↓          → valuable (采集)
IF r̄↓ AND ρ↑ AND D↑                  → redundant (压缩)
IF r̄↑ AND ρ↓ AND C↓                  → noisy (降权/丢弃)
ELSE                                  → valuable (保守)
\end{verbatim}

代码：\texttt{valuation/classifier.py} 的
\texttt{DataClassifier.classify\_state()}

\paragraph{两层描述符 (Proposition
6)}\label{ux4e24ux5c42ux63cfux8ff0ux7b26-proposition-6}

\begin{verbatim}
Layer 1: φ_human(x) → KMeans → coarse states {S1, S2, ..., SK}
Layer 2: ∀ S_k:
   1. 计算 MI(feature_dim, error) for each feature dim
   2. 选 Top-d 误差相关维度
   3. 仅在误差子空间聚类 → error_states
   4. 输出两层级联划分
\end{verbatim}

代码：\texttt{state/two\_layer.py} 的
\texttt{TwoLayerStateDiscovery.discover()}

\begin{center}\rule{0.5\linewidth}{0.5pt}\end{center}

\subsection{3. 接口层（Encoder
抽象基类）------连接基底与应用}\label{ux63a5ux53e3ux5c42encoder-ux62bdux8c61ux57faux7c7bux8fdeux63a5ux57faux5e95ux4e0eux5e94ux7528}

\subsubsection{3.1
SCXStateEncoder（核心接口）}\label{scxstateencoderux6838ux5fc3ux63a5ux53e3}

\begin{Shaded}
\begin{Highlighting}[]
\KeywordTok{class}\NormalTok{ SCXStateEncoder(ABC):}
    \AttributeTok{@abstractmethod}
    \KeywordTok{def}\NormalTok{ encode(}\VariableTok{self}\NormalTok{, x: Any) }\OperatorTok{{-}\textgreater{}}\NormalTok{ np.ndarray:}
        \CommentTok{"""x → φ(x) ∈ R\^{}d"""}

    \AttributeTok{@abstractmethod}
    \KeywordTok{def}\NormalTok{ distance(}\VariableTok{self}\NormalTok{, a: np.ndarray, b: np.ndarray) }\OperatorTok{{-}\textgreater{}} \BuiltInTok{float}\NormalTok{:}
        \CommentTok{"""特征空间中的距离"""}

    \AttributeTok{@abstractmethod}
    \KeywordTok{def}\NormalTok{ cluster(}\VariableTok{self}\NormalTok{, X: np.ndarray, n\_clusters: }\BuiltInTok{int}\NormalTok{, }\OperatorTok{**}\NormalTok{kw) }\OperatorTok{{-}\textgreater{}} \BuiltInTok{tuple}\NormalTok{[labels, centroids]:}
        \CommentTok{"""聚类划分状态"""}

    \KeywordTok{def}\NormalTok{ batch\_encode(}\VariableTok{self}\NormalTok{, X: }\BuiltInTok{list}\NormalTok{) }\OperatorTok{{-}\textgreater{}}\NormalTok{ np.ndarray:}
        \CommentTok{"""批量编码（可重写做并行）"""}

    \KeywordTok{def}\NormalTok{ get\_default\_config(}\VariableTok{self}\NormalTok{) }\OperatorTok{{-}\textgreater{}} \BuiltInTok{dict}\NormalTok{:}
        \CommentTok{"""默认状态发现配置 \{n\_states, cluster\_method\}"""}

    \KeywordTok{def}\NormalTok{ get\_feature\_dim(}\VariableTok{self}\NormalTok{) }\OperatorTok{{-}\textgreater{}} \BuiltInTok{int}\NormalTok{:}
        \CommentTok{"""返回特征维度 d"""}
\end{Highlighting}
\end{Shaded}

\textbf{这就是唯一必须实现的接口。} 无论你做 MLIP、CNN、Transformer、LLM
embedding------只要实现这 3 个抽象方法，SCX 的整个基底就能跑。

\subsubsection{3.2
SCXExpert（轻量专家接口）}\label{scxexpertux8f7bux91cfux4e13ux5bb6ux63a5ux53e3}

\begin{Shaded}
\begin{Highlighting}[]
\KeywordTok{class}\NormalTok{ SCXExpert(ABC):}
    \AttributeTok{@abstractmethod}
    \KeywordTok{def}\NormalTok{ predict(}\VariableTok{self}\NormalTok{, x: Any) }\OperatorTok{{-}\textgreater{}}\NormalTok{ Any:}
        \CommentTok{"""专家预测"""}

    \KeywordTok{def}\NormalTok{ cost(}\VariableTok{self}\NormalTok{) }\OperatorTok{{-}\textgreater{}} \BuiltInTok{float}\NormalTok{:}
        \CommentTok{"""调用成本 (默认 1.0)"""}
\end{Highlighting}
\end{Shaded}

\begin{center}\rule{0.5\linewidth}{0.5pt}\end{center}

\subsection{4. 领域适配层（Encoder 实现 + YAML
配置）}\label{ux9886ux57dfux9002ux914dux5c42encoder-ux5b9eux73b0-yaml-ux914dux7f6e}

\subsubsection{4.1 已实现的
Encoder}\label{ux5df2ux5b9eux73b0ux7684-encoder}

\begin{longtable}[]{@{}
  >{\raggedright\arraybackslash}p{(\linewidth - 8\tabcolsep) * \real{0.1875}}
  >{\raggedright\arraybackslash}p{(\linewidth - 8\tabcolsep) * \real{0.1250}}
  >{\raggedright\arraybackslash}p{(\linewidth - 8\tabcolsep) * \real{0.3750}}
  >{\raggedright\arraybackslash}p{(\linewidth - 8\tabcolsep) * \real{0.1875}}
  >{\raggedright\arraybackslash}p{(\linewidth - 8\tabcolsep) * \real{0.1250}}@{}}
\toprule\noalign{}
\begin{minipage}[b]{\linewidth}\raggedright
Encoder
\end{minipage} & \begin{minipage}[b]{\linewidth}\raggedright
文件
\end{minipage} & \begin{minipage}[b]{\linewidth}\raggedright
用于什么网络/模型
\end{minipage} & \begin{minipage}[b]{\linewidth}\raggedright
特征维度
\end{minipage} & \begin{minipage}[b]{\linewidth}\raggedright
领域
\end{minipage} \\
\midrule\noalign{}
\endhead
\bottomrule\noalign{}
\endlastfoot
\texttt{MLIPEncoder} & \texttt{encoders/mlip.py} & ACE/NEP/MACE
势函数（手工描述符：配位数/键长/体积\ldots） & 12-dim & 材料/DFT \\
\texttt{VisionEncoder} & \texttt{encoders/vision.py} & SimpleCNN /
ResNet18 / ResNet50 & 256-dim (simple\_cnn) & 图像分类 \\
\texttt{TabularEncoder} & \texttt{encoders/tabular.py} &
任何表格模型（归一化+one-hot） & 可变 & 表格数据 \\
\texttt{ErrorDrivenEncoder} & \texttt{encoders/error\_driven.py} &
\textbf{Layer 2 专用}------包装任意 Layer 1 encoder &
可变（误差相关子空间） & 全部 \\
\end{longtable}

\subsubsection{4.2 如何为新领域写 Encoder（3
步）}\label{ux5982ux4f55ux4e3aux65b0ux9886ux57dfux5199-encoder3-ux6b65}

\begin{Shaded}
\begin{Highlighting}[]
\CommentTok{\# 例：为 LLM 文本分类写 Encoder}
\ImportTok{from}\NormalTok{ scx.encoders.base }\ImportTok{import}\NormalTok{ SCXStateEncoder}
\ImportTok{import}\NormalTok{ numpy }\ImportTok{as}\NormalTok{ np}

\KeywordTok{class}\NormalTok{ LLMEncoder(SCXStateEncoder):}
    \KeywordTok{def} \FunctionTok{\_\_init\_\_}\NormalTok{(}\VariableTok{self}\NormalTok{, model\_name}\OperatorTok{=}\StringTok{"bert{-}base"}\NormalTok{):}
        \VariableTok{self}\NormalTok{.model }\OperatorTok{=}\NormalTok{ load\_embedding\_model(model\_name)  }\CommentTok{\# 你的 embedding 模型}
        \VariableTok{self}\NormalTok{.\_feature\_dim }\OperatorTok{=} \DecValTok{768}

    \KeywordTok{def}\NormalTok{ encode(}\VariableTok{self}\NormalTok{, text: }\BuiltInTok{str}\NormalTok{) }\OperatorTok{{-}\textgreater{}}\NormalTok{ np.ndarray:}
        \ControlFlowTok{return} \VariableTok{self}\NormalTok{.model.embed(text)  }\CommentTok{\# text → 768{-}dim vector}

    \KeywordTok{def}\NormalTok{ distance(}\VariableTok{self}\NormalTok{, a, b):}
        \ControlFlowTok{return} \BuiltInTok{float}\NormalTok{(np.linalg.norm(a }\OperatorTok{{-}}\NormalTok{ b))  }\CommentTok{\# 欧氏距离}

    \KeywordTok{def}\NormalTok{ cluster(}\VariableTok{self}\NormalTok{, X, n\_clusters}\OperatorTok{=}\DecValTok{10}\NormalTok{, }\OperatorTok{**}\NormalTok{kw):}
        \ImportTok{from}\NormalTok{ sklearn.cluster }\ImportTok{import}\NormalTok{ KMeans}
\NormalTok{        km }\OperatorTok{=}\NormalTok{ KMeans(n\_clusters}\OperatorTok{=}\NormalTok{n\_clusters, random\_state}\OperatorTok{=}\DecValTok{42}\NormalTok{)}
        \ControlFlowTok{return}\NormalTok{ km.fit\_predict(X), km.cluster\_centers\_}
\end{Highlighting}
\end{Shaded}

\textbf{3 步接入新领域}： 1. 实现 \texttt{SCXStateEncoder}（如上） 2. 写
YAML 配置（\texttt{domains/llm.yaml}） 3.
\texttt{DomainRegistry.register("llm",\ "domains/llm.yaml")}

\subsubsection{4.3
已注册的领域配置}\label{ux5df2ux6ce8ux518cux7684ux9886ux57dfux914dux7f6e}

\begin{longtable}[]{@{}
  >{\raggedright\arraybackslash}p{(\linewidth - 8\tabcolsep) * \real{0.1500}}
  >{\raggedright\arraybackslash}p{(\linewidth - 8\tabcolsep) * \real{0.2500}}
  >{\raggedright\arraybackslash}p{(\linewidth - 8\tabcolsep) * \real{0.2250}}
  >{\raggedright\arraybackslash}p{(\linewidth - 8\tabcolsep) * \real{0.1500}}
  >{\raggedright\arraybackslash}p{(\linewidth - 8\tabcolsep) * \real{0.2250}}@{}}
\toprule\noalign{}
\begin{minipage}[b]{\linewidth}\raggedright
领域
\end{minipage} & \begin{minipage}[b]{\linewidth}\raggedright
YAML 文件
\end{minipage} & \begin{minipage}[b]{\linewidth}\raggedright
Encoder
\end{minipage} & \begin{minipage}[b]{\linewidth}\raggedright
专家
\end{minipage} & \begin{minipage}[b]{\linewidth}\raggedright
阈值特色
\end{minipage} \\
\midrule\noalign{}
\endhead
\bottomrule\noalign{}
\endlastfoot
\textbf{mlip} & \texttt{domains/mlip.yaml} & \texttt{MLIPEncoder}
(12-dim ACE描述符) & ACE, NEP, MACE, DFT & error\_high=0.05 eV/atom \\
\textbf{cifar\_image} & \texttt{domains/cifar.yaml} &
\texttt{VisionEncoder} (simple\_cnn) & cnn\_small, cnn\_large, resnet50
& error\_high=0.3 (loss) \\
\textbf{health\_image} & \texttt{domains/health.yaml} &
\texttt{VisionEncoder} (simple\_cnn, 224x224) & ai\_v1, ai\_v2,
junior\_doctor, senior\_doctor & expert\_gap=0.25（人机混合） \\
\end{longtable}

\begin{center}\rule{0.5\linewidth}{0.5pt}\end{center}

\subsection{5.
完整脚本作用清单}\label{ux5b8cux6574ux811aux672cux4f5cux7528ux6e05ux5355}

\subsubsection{\texorpdfstring{5.1
基底脚本（\texttt{src/scx/}）}{5.1 基底脚本（src/scx/）}}\label{ux57faux5e95ux811aux672csrcscx}

\begin{verbatim}
src/scx/
├── __init__.py                    导出 SCXFramework, SCXConfig, 所有 Encoder, DomainRegistry
├── core/
│   ├── config.py                  SCXConfig — 全局配置 (K, thresholds, budget)
│   ├── framework.py               SCXFramework — 总调度器 fit/compress/value/save
│   ├── metrics.py                 SCXMetrics — 指标追踪
│   └── online.py                  OnlineSCXFramework — 流式/在线版本
├── encoders/
│   ├── base.py                    SCXStateEncoder, SCXExpert — 抽象基类（唯一接口）
│   ├── mlip.py                    MLIPEncoder — ACE/势函数 (12-dim 手工描述符)
│   ├── vision.py                  VisionEncoder — 图像 (CNN/ResNet backbone)
│   ├── tabular.py                 TabularEncoder — 表格数据
│   └── error_driven.py            ErrorDrivenEncoder — Layer 2 错误驱动聚类
├── domains/
│   ├── registry.py                DomainRegistry — 注册/查找/加载 YAML 配置
│   ├── mlip.yaml                  MLIP 领域配置
│   ├── cifar.yaml                 CIFAR 领域配置
│   └── health.yaml                Medical 领域配置
├── state/
│   ├── discovery.py               StateDiscovery — 聚类
│   ├── two_layer.py               TwoLayerStateDiscovery — 两层描述符管线
│   ├── assignment.py              StateAssignment — 新样本映射
│   ├── space.py                   StateSpace — 状态空间数据结构
│   ├── metrics.py                 状态间距离/相似度
│   └── robustness.py              StateRobustness — 鲁棒性分析
├── valuation/
│   ├── classifier.py              DataClassifier — 四分类器
│   ├── state_value.py             StateValue — V(s) [deprecated] + theorem-based methods
│   ├── noise_score.py             NoiseScore — 噪声评分
│   ├── redundancy.py              RedundancyScore — 冗余评分
│   ├── learnability.py            LearnabilityScore — 可学习性
│   ├── adaptive.py                AdaptiveThreshold — 自适应阈值
│   └── influence.py               StateConditionedInfluence — 影响力
├── yajie.py                      YajieCleaner — 雅洁数据清理器（定理驱动去噪）
├── expert/
│   ├── registry.py                ExpertRegistry — 专家注册
│   ├── reliability.py             ExpertReliability — SCX_m(s) 计算
│   ├── router.py                  ExpertRouter — hard/soft 路由
│   └── conflict.py                专家冲突仲裁
├── action/
│   ├── compress.py                CompressStrategy — 压缩
│   ├── acquisition.py             主动采集
│   └── policy.py                  ActionPolicy — 决策策略
└── utils/
    ├── data_loader.py             数据加载
    ├── evaluation.py              评估工具
    ├── visualization.py           可视化
    └── helpers.py                 辅助函数
\end{verbatim}

\subsubsection{\texorpdfstring{5.2
实验脚本（\texttt{experiments/}）}{5.2 实验脚本（experiments/）}}\label{ux5b9eux9a8cux811aux672cexperiments}

\begin{longtable}[]{@{}
  >{\raggedright\arraybackslash}p{(\linewidth - 4\tabcolsep) * \real{0.2143}}
  >{\raggedright\arraybackslash}p{(\linewidth - 4\tabcolsep) * \real{0.2143}}
  >{\raggedright\arraybackslash}p{(\linewidth - 4\tabcolsep) * \real{0.5714}}@{}}
\toprule\noalign{}
\begin{minipage}[b]{\linewidth}\raggedright
脚本
\end{minipage} & \begin{minipage}[b]{\linewidth}\raggedright
作用
\end{minipage} & \begin{minipage}[b]{\linewidth}\raggedright
调用的 SCX 模块
\end{minipage} \\
\midrule\noalign{}
\endhead
\bottomrule\noalign{}
\endlastfoot
\texttt{mlip\_case/run\_scx\_on\_aln\_v3.py} & \textbf{一层} SCX 分析
AlN v3（534帧） & MLIPEncoder, StateDiscovery, DataClassifier,
NoiseScore, StateValue \\
\texttt{mlip\_case/run\_scx\_two\_layer.py} & \textbf{两层} SCX 分析 AlN
v3 & MLIPEncoder + ErrorDrivenEncoder + TwoLayerStateDiscovery \\
\texttt{mlip\_case/compare\_analysis.py} & SCX vs DFT 对比 & 读取 v3
训练结果 + SCX 结果 \\
\texttt{mlip\_case/analyze\_overlap.py} & 噪声帧重叠分析 & 交叉对比 SCX
噪声标记 vs REPORT 最差帧 \\
\texttt{mlip\_case/final\_check.py} & 最终验证统计 & Pearson r, RMSE
预估值 \\
\texttt{cifar/run\_cifar10\_noise.py} & CIFAR-10 噪声检测实验 &
VisionEncoder, NoiseScore \\
\texttt{cifar/run\_cifar10\_compress.py} & CIFAR-10 压缩实验 &
VisionEncoder, CompressStrategy \\
\texttt{cifar/run\_cifar100\_routing.py} & CIFAR-100 路由实验 &
VisionEncoder, ExpertRouter \\
\texttt{cifar/run\_baselines.py} & CIFAR baseline 对比 & DataClassifier
vs Random/LOO/Shapley \\
\texttt{synthetic/data\_generator.py} & 合成数据生成 & 生成 2D
状态条件数据集 \\
\texttt{synthetic/run\_experiment.py} & 合成数据验证 & 全管线 \\
(待实现) \texttt{distill/devirus.py} & \textbf{蒸馏去病毒化} & Teacher
势函数蒸馏数据 + SCX 两层过滤 \\
\end{longtable}

\subsubsection{\texorpdfstring{5.3 scx-health
实验脚本（\texttt{scx-health/experiments/}）}{5.3 scx-health 实验脚本（scx-health/experiments/）}}\label{scx-health-ux5b9eux9a8cux811aux672cscx-healthexperiments}

\begin{longtable}[]{@{}
  >{\raggedright\arraybackslash}p{(\linewidth - 4\tabcolsep) * \real{0.3000}}
  >{\raggedright\arraybackslash}p{(\linewidth - 4\tabcolsep) * \real{0.3000}}
  >{\raggedright\arraybackslash}p{(\linewidth - 4\tabcolsep) * \real{0.4000}}@{}}
\toprule\noalign{}
\begin{minipage}[b]{\linewidth}\raggedright
脚本
\end{minipage} & \begin{minipage}[b]{\linewidth}\raggedright
作用
\end{minipage} & \begin{minipage}[b]{\linewidth}\raggedright
数据集
\end{minipage} \\
\midrule\noalign{}
\endhead
\bottomrule\noalign{}
\endlastfoot
\texttt{compress/run\_compress.py} & SCX-Compress 完整版 (含 Coreset) &
PathMNIST 90K \\
\texttt{compress/run\_compress\_v2\_quick.py} & SCX-Compress 快速版 (无
Coreset) & PathMNIST 90K \\
\texttt{noise/run\_noise.py} & SCX-Noise 均匀+按类别噪声 & DermaMNIST
7K \\
\texttt{routing/run\_routing.py} & SCX-Routing 多专家路由 & BloodMNIST
12K \\
\end{longtable}

\subsubsection{\texorpdfstring{5.4
测试脚本（\texttt{tests/}）}{5.4 测试脚本（tests/）}}\label{ux6d4bux8bd5ux811aux672ctests}

\begin{longtable}[]{@{}
  >{\raggedright\arraybackslash}p{(\linewidth - 4\tabcolsep) * \real{0.2609}}
  >{\raggedright\arraybackslash}p{(\linewidth - 4\tabcolsep) * \real{0.3913}}
  >{\raggedright\arraybackslash}p{(\linewidth - 4\tabcolsep) * \real{0.3478}}@{}}
\toprule\noalign{}
\begin{minipage}[b]{\linewidth}\raggedright
文件
\end{minipage} & \begin{minipage}[b]{\linewidth}\raggedright
覆盖模块
\end{minipage} & \begin{minipage}[b]{\linewidth}\raggedright
测试数
\end{minipage} \\
\midrule\noalign{}
\endhead
\bottomrule\noalign{}
\endlastfoot
\texttt{test\_state.py} & StateDiscovery, StateSpace, StateAssignment,
两层 & \textasciitilde100 \\
\texttt{test\_valuation.py} & DataClassifier, NoiseScore, StateValue,
Learnability, Redundancy & \textasciitilde80 \\
\texttt{test\_expert.py} & ExpertRouter, ExpertReliability &
\textasciitilde60 \\
\texttt{test\_expert\_module.py} & ExpertRegistry & \textasciitilde30 \\
\texttt{test\_action.py} & CompressStrategy, ActionPolicy &
\textasciitilde50 \\
\texttt{test\_two\_layer.py} & TwoLayerStateDiscovery,
ErrorDrivenEncoder & \textasciitilde30 \\
\texttt{test\_robustness.py} & StateRobustness & \textasciitilde20 \\
\texttt{test\_online.py} & OnlineSCXFramework & \textasciitilde20 \\
\texttt{test\_utils.py} & 辅助函数 & \textasciitilde10 \\
\texttt{test\_yajie.py} & YajieCleaner & \textasciitilde3 \\
\texttt{conftest.py} & 共享 fixtures & --- \\
\end{longtable}

\textbf{总计 \textasciitilde427 tests}

\begin{center}\rule{0.5\linewidth}{0.5pt}\end{center}

\subsection{6.
不同应用场景的调用方式}\label{ux4e0dux540cux5e94ux7528ux573aux666fux7684ux8c03ux7528ux65b9ux5f0f}

\subsubsection{6.1 MLIP 场景（AlN v3）}\label{mlip-ux573aux666faln-v3}

\begin{Shaded}
\begin{Highlighting}[]
\ImportTok{from}\NormalTok{ scx.encoders.mlip }\ImportTok{import}\NormalTok{ MLIPEncoder}
\ImportTok{from}\NormalTok{ scx.state.discovery }\ImportTok{import}\NormalTok{ StateDiscovery}
\ImportTok{from}\NormalTok{ scx.valuation.classifier }\ImportTok{import}\NormalTok{ DataClassifier}
\ImportTok{from}\NormalTok{ scx.valuation.noise\_score }\ImportTok{import}\NormalTok{ NoiseScore}

\CommentTok{\# 1. 编码}
\NormalTok{encoder }\OperatorTok{=}\NormalTok{ MLIPEncoder(species}\OperatorTok{=}\NormalTok{[}\StringTok{"Al"}\NormalTok{, }\StringTok{"N"}\NormalTok{])}
\NormalTok{features }\OperatorTok{=}\NormalTok{ encoder.batch\_encode(atoms\_list)  }\CommentTok{\# 534 x 12}

\CommentTok{\# 2. 状态发现}
\NormalTok{discovery }\OperatorTok{=}\NormalTok{ StateDiscovery(method}\OperatorTok{=}\StringTok{"kmeans"}\NormalTok{, n\_states}\OperatorTok{=}\DecValTok{20}\NormalTok{)}
\NormalTok{labels }\OperatorTok{=}\NormalTok{ discovery.fit\_predict(features)}

\CommentTok{\# 3. 计算残差 (fmax)}
\NormalTok{fmax }\OperatorTok{=}\NormalTok{ compute\_fmax(atoms\_list)}

\CommentTok{\# 4. 四分类}
\NormalTok{classifier }\OperatorTok{=}\NormalTok{ DataClassifier(\{}\StringTok{"error\_high"}\NormalTok{: }\FloatTok{3.0}\NormalTok{\})  }\CommentTok{\# MLIP 专用阈值}
\ControlFlowTok{for}\NormalTok{ s }\KeywordTok{in}\NormalTok{ states:}
\NormalTok{    cat }\OperatorTok{=}\NormalTok{ classifier.classify\_state(r\_bar, prop, consistency, redundancy, noise)}

\CommentTok{\# 5. 两层方法}
\ImportTok{from}\NormalTok{ scx.state.two\_layer }\ImportTok{import}\NormalTok{ TwoLayerStateDiscovery}
\NormalTok{two\_layer }\OperatorTok{=}\NormalTok{ TwoLayerStateDiscovery(encoder)}
\NormalTok{result }\OperatorTok{=}\NormalTok{ two\_layer.discover(atoms\_list, fmax, layer1\_k}\OperatorTok{=}\DecValTok{20}\NormalTok{, layer2\_k}\OperatorTok{=}\DecValTok{5}\NormalTok{)}
\end{Highlighting}
\end{Shaded}

\subsubsection{6.2
图像场景（CIFAR/MedMNIST）}\label{ux56feux50cfux573aux666fcifarmedmnist}

\begin{Shaded}
\begin{Highlighting}[]
\ImportTok{from}\NormalTok{ scx.encoders.vision }\ImportTok{import}\NormalTok{ VisionEncoder}

\CommentTok{\# 编码器换成 VisionEncoder，其他一样}
\NormalTok{encoder }\OperatorTok{=}\NormalTok{ VisionEncoder(backbone}\OperatorTok{=}\StringTok{"simple\_cnn"}\NormalTok{)}
\NormalTok{features }\OperatorTok{=}\NormalTok{ encoder.batch\_encode(images)  }\CommentTok{\# N x 256}
\CommentTok{\# 后续管线完全不变}
\end{Highlighting}
\end{Shaded}

\subsubsection{6.3 表格场景}\label{ux8868ux683cux573aux666f}

\begin{Shaded}
\begin{Highlighting}[]
\ImportTok{from}\NormalTok{ scx.encoders.tabular }\ImportTok{import}\NormalTok{ TabularEncoder}

\NormalTok{encoder }\OperatorTok{=}\NormalTok{ TabularEncoder(normalize}\OperatorTok{=}\VariableTok{True}\NormalTok{)}
\NormalTok{features }\OperatorTok{=}\NormalTok{ encoder.batch\_encode(dataframe)  }\CommentTok{\# N x d}
\CommentTok{\# 后续管线完全不变}
\end{Highlighting}
\end{Shaded}

\subsubsection{6.4 LLM
场景（待实现）}\label{llm-ux573aux666fux5f85ux5b9eux73b0}

\begin{Shaded}
\begin{Highlighting}[]
\CommentTok{\# 只需写一个 LLMEncoder，其他全不变}
\ImportTok{from}\NormalTok{ scx.encoders.base }\ImportTok{import}\NormalTok{ SCXStateEncoder}

\KeywordTok{class}\NormalTok{ LLMEncoder(SCXStateEncoder):}
    \KeywordTok{def}\NormalTok{ encode(}\VariableTok{self}\NormalTok{, text): }\ControlFlowTok{return}\NormalTok{ embedding\_model(text)}
    \KeywordTok{def}\NormalTok{ distance(}\VariableTok{self}\NormalTok{, a, b): }\ControlFlowTok{return}\NormalTok{ cosine\_distance(a, b)}
    \KeywordTok{def}\NormalTok{ cluster(}\VariableTok{self}\NormalTok{, X, n, }\OperatorTok{**}\NormalTok{kw): }\ControlFlowTok{return}\NormalTok{ KMeans(n).fit\_predict(X), centroids}

\CommentTok{\# 后续管线完全不变}
\end{Highlighting}
\end{Shaded}

\begin{center}\rule{0.5\linewidth}{0.5pt}\end{center}

\subsection{7. 数据流图（一次完整的 SCX
调用）}\label{ux6570ux636eux6d41ux56feux4e00ux6b21ux5b8cux6574ux7684-scx-ux8c03ux7528}

\begin{verbatim}
输入数据 X (原子构型/图像/文本/表格行)
    │
    ▼
┌──────────────────────┐
│ ① Encoder.encode()  │  ← 领域唯一差异点
│   X → φ(X) ∈ R^d    │
└──────────┬───────────┘
           │ φ(X)  numpy array (N, d)
           ▼
┌──────────────────────┐
│ ② StateDiscovery    │
│   cluster(φ(X), K)  │
│   → labels, centroids│
└──────────┬───────────┘
           │ state_labels (N,)
           ▼
┌──────────────────────┐
│ ③ 残差计算          │
│   r_i = ℓ(f(x_i), y_i)│ ← 需要模型预测+标签
└──────────┬───────────┘
           │ residuals (N,)
           ▼
┌──────────────────────┐
│ ④ 四维估值 (per state)│
│   r̄(s) ρ(s) L(s) D(s)│
│   NoiseScore(s)      │
│   SCX_m(s) for each m│
└──────────┬───────────┘
           │ state_metrics dict
           ▼
┌──────────────────────┐
│ ⑤ DataClassifier    │
│   classify_state()   │
│   → Valuable/Redun/  │
│     Noisy/ExpertDep  │
└──────────┬───────────┘
           │ category per state
           ▼
┌──────────────────────┐
│ ⑥ ActionPolicy      │
│   → acquire/compress │
│     /downweight/route│
└──────────┬───────────┘
           │
           ▼
┌──────────────────────┐
│ ⑦ 反馈 (闭环)       │
│   真实结果回传        │
│   重校准阈值          │
└──────────────────────┘
\end{verbatim}

\textbf{第①步是唯一因领域而异的部分。} ②-⑦ 对所有领域完全相同。

\begin{center}\rule{0.5\linewidth}{0.5pt}\end{center}

\subsection{8.
两类两层的实现对照}\label{ux4e24ux7c7bux4e24ux5c42ux7684ux5b9eux73b0ux5bf9ux7167}

\begin{longtable}[]{@{}
  >{\raggedright\arraybackslash}p{(\linewidth - 4\tabcolsep) * \real{0.3333}}
  >{\raggedright\arraybackslash}p{(\linewidth - 4\tabcolsep) * \real{0.3333}}
  >{\raggedright\arraybackslash}p{(\linewidth - 4\tabcolsep) * \real{0.3333}}@{}}
\toprule\noalign{}
\begin{minipage}[b]{\linewidth}\raggedright
\end{minipage} & \begin{minipage}[b]{\linewidth}\raggedright
EGP 两层 (V3 Phase C)
\end{minipage} & \begin{minipage}[b]{\linewidth}\raggedright
SCX 两层 (v4.0)
\end{minipage} \\
\midrule\noalign{}
\endhead
\bottomrule\noalign{}
\endlastfoot
\textbf{目的} & 合并势函数时消除参数冗余 & 发现噪声/冗余/高价值状态 \\
\textbf{Layer 1} & shared c\_0 (共用) & 人工描述符粗聚类 \\
\textbf{Layer 2} & element correction c\_Z &
ErrorDrivenEncoder（误差子空间聚类） \\
\textbf{操作对象} & 模型系数 & 训练数据 \\
\textbf{关键操作} & post-hoc projection (violation→0) & 互信息筛选 +
误差子空间聚类 \\
\textbf{实物产出} & gauge-fixed 势函数 & 数据四分类 + 去噪建议 \\
\end{longtable}

\begin{center}\rule{0.5\linewidth}{0.5pt}\end{center}

\subsection{9.
关键对照表：什么场景用什么}\label{ux5173ux952eux5bf9ux7167ux8868ux4ec0ux4e48ux573aux666fux7528ux4ec0ux4e48}

\begin{longtable}[]{@{}
  >{\raggedright\arraybackslash}p{(\linewidth - 6\tabcolsep) * \real{0.2157}}
  >{\raggedright\arraybackslash}p{(\linewidth - 6\tabcolsep) * \real{0.2549}}
  >{\raggedright\arraybackslash}p{(\linewidth - 6\tabcolsep) * \real{0.2157}}
  >{\raggedright\arraybackslash}p{(\linewidth - 6\tabcolsep) * \real{0.3137}}@{}}
\toprule\noalign{}
\begin{minipage}[b]{\linewidth}\raggedright
你要做的事
\end{minipage} & \begin{minipage}[b]{\linewidth}\raggedright
用的 Encoder
\end{minipage} & \begin{minipage}[b]{\linewidth}\raggedright
用的 YAML
\end{minipage} & \begin{minipage}[b]{\linewidth}\raggedright
用的关键基底模块
\end{minipage} \\
\midrule\noalign{}
\endhead
\bottomrule\noalign{}
\endlastfoot
分析 AlN DFT 训练数据质量 & \texttt{MLIPEncoder} & \texttt{mlip.yaml} &
TwoLayerDiscovery, DataClassifier, NoiseScore \\
CIFAR 图像去噪/压缩 & \texttt{VisionEncoder} (simple\_cnn) &
\texttt{cifar.yaml} & NoiseScore, CompressStrategy \\
MedMNIST 医学图像路由 & \texttt{VisionEncoder} (simple\_cnn) &
\texttt{health.yaml} & ExpertRouter, DataClassifier \\
表格数据估值 (UCI) & \texttt{TabularEncoder} & (新建 tabular.yaml) &
StateValue, DataClassifier \\
LLM embedding 分析 & (待写 \texttt{LLMEncoder}) & (新建 llm.yaml) &
全部通用 \\
FEM 仿真参数空间 & (待写 \texttt{FEMEncoder}) & (新建 fem.yaml) &
TwoLayerDiscovery, StateValue \\
在线/流式数据 & 任意 Encoder & 任意 YAML & \texttt{OnlineSCXFramework},
\texttt{OnlineStateTracker} \\
\end{longtable}

\begin{center}\rule{0.5\linewidth}{0.5pt}\end{center}

\subsection{10. 一句话总结}\label{ux4e00ux53e5ux8bddux603bux7ed3}

\begin{quote}
\textbf{Encoder 是唯一的领域插头，基底是不变的核心引擎。}\\
MLIP 用 \texttt{MLIPEncoder}（手工描述符），图像用
\texttt{VisionEncoder}（CNN），LLM 你写个 \texttt{LLMEncoder}。\\
插头一换，整个 SCX
管线就能在新领域跑起来------StateDiscovery、DataClassifier、NoiseScore、ExpertRouter
全部复用。
\end{quote}

\end{document}
